\documentclass[11pt,a4paper]{article}
\usepackage[T1]{fontenc}
\usepackage[utf8]{inputenc}
\usepackage{lmodern}
\usepackage{amsmath,amssymb}
\usepackage[margin=25mm]{geometry}
\usepackage{longtable,booktabs,array,calc}
\usepackage[hidelinks]{hyperref}
\hypersetup{pdftitle={Inserting a point, subdividing a cell: how the limit is built},pdfauthor={navstokgap research working notes}}
\newcommand{\tightlist}{\setlength{\itemsep}{0pt}\setlength{\parskip}{0pt}}
\setlength{\emergencystretch}{3em}
\setcounter{secnumdepth}{0}
\begin{document}
\hypertarget{inserting-a-point-subdividing-a-cell-how-the-limit-is-built}{%
\section{Inserting a point, subdividing a cell: how the limit is
built}\label{inserting-a-point-subdividing-a-cell-how-the-limit-is-built}}

\textbf{Result and direction, 2026-09-26.} The user's insertion
\((t_0,t_1,t_2,t_3)\mapsto(t_0,t_1,t_2,t_{2.6},t_3)\) identifies the
right local question: after adding a variable, how are the old
observations and the dynamics recovered? For Galileo's constant force,
eliminating the new point gives an exact composition law and an explicit
action correction. A summable-defect lemma then constructs a limit
independent of the sequence of insertions. In \(1+1\)-dimensional pure
Yang--Mills, a heat-kernel convolution gives exact consistency under
inserting an edge. In \(1+2\) and \(1+3\) dimensions the corresponding
integration produces interactions among several boundary variables;
their control is the renormalization problem.

This is a constructive development of the
\href{three-continuum-limits.md}{joint-paper plan}. The elementary
proofs below assemble established ingredients, with no novelty claim.
They produce the constant-force limit and the exact two-dimensional
comparison, and specify sufficient estimates for further limits. The
positive action necessity and the four-dimensional \(SU(3)\) gap remain
the two main goals; the pion remains a symmetry benchmark.

\textbf{Foundational purpose (user, 2026-09-30).} A reconstruction of
quantum theory from the limits of classical mechanics is judged here by
the constructions and estimates it enables, especially for continuum
existence and the physical gap. Explaining why a change of resolution
introduces an RG transformation would already be useful. The exact
blocking identities below provide that structural explanation; the
quantum weighting and its positive action scale remain stated premises.

\textbf{All refinements and scale normalization, 2026-09-30 (GPT-6.1
Sol; mathematical bounds in §§3c--3d checked by GPT-6 Astra).} The
harmonic curve is constructed without assuming it, uniformly over all
admitted partitions and arbitrary mixtures of start and end kicks. The
two finite conventions differ, while their limits agree. This gives a
precise comparison with the user's RG/'t Hooft thread. Mass cancellation
in classical motion, classical gauge scale invariance and quantum
dimensional transmutation are separated in §3d: refinement independence
of a curve leaves its action normalization undetermined.

\textbf{Further result, 2026-09-29 (GPT-6.1 Sol; written, §3b checked by
GPT-6 Astra).} For the end-kick harmonic polygon, §3b proves a finite
Cauchy estimate uniform over unequal partitions and observation times. A
weighted absolute-value norm supplies stability and telescoping
insertion bounds without assuming the limiting curve. A finite
force-balance residual tends to zero with the same mesh. This develops
the specific norm obstacle exposed by the sister Lean formalization; it
supplies no action floor or measurement law.

\hypertarget{what-an-insertion-has-to-preserve}{%
\subsection{1. What an insertion has to
preserve}\label{what-an-insertion-has-to-preserve}}

Three operations occur at a new time \(r\) between \(s\) and \(t\).

\begin{enumerate}
\def\labelenumi{\arabic{enumi}.}
\tightlist
\item
  A classical calculation adds \(q(r)\) and eliminates it by the
  equations of motion, or by stationary action with fixed endpoints.
\item
  A quantum propagator integrates over an unobserved intermediate
  coordinate: \(K(t,s)=\int K(t,r)K(r,s)\,dq(r)\).
\item
  A physical mark introduces an apparatus and a record. Forgetting that
  record applies the mark's nonselective channel to the body.
\end{enumerate}

Each is a precise composition question. In the third case the
nonselective channel must act as the identity on retained observables if
one is to recover the unmarked dynamics. Trace preservation alone does
not establish that. Thus inserting a variable in a calculation and
inserting an actual measurement have separate refinement laws.

For a gauge lattice, multiplying fine links along an old edge defines
the old holonomy, while integrating the other variables defines its new
probability law. Both the projection of variables and the law must be
specified. A choice of fine variables over a coarse configuration also
needs a conditional distribution; coarse data alone do not choose that
distribution.

\hypertarget{newtons-one-point-insertion-with-all-constants}{%
\subsection{2. Newton's one-point insertion, with all
constants}\label{newtons-one-point-insertion-with-all-constants}}

Take transverse mass \(M>0\), force \(F\), and a cell of duration
\(h=u+v\), split after \(u>0\), with \(v>0\). This is the symmetric
impulsive construction of the \href{polygon-lift-phase.md}{polygon
note}, a modern constant-force realization of the inscribed-chord
comparison. It does not assert that every historical force polygon used
this update.

Write a kick and a free drift as

\[K_J(q,p)=(q,p+J),\qquad D_h(q,p)=(q+hp/M,p).\]

The cell map is

\[S_h=K_{Fh/2}D_hK_{Fh/2},\qquad
S_h(q,p)=\left(q+\frac{hp}{M}+\frac{Fh^2}{2M},p+Fh\right).\]

The input momentum is before the first half-kick and the output is after
the last. At a shared vertex the retained momentum is the one between
the two adjacent half-kicks; the incoming and outgoing chord momenta
depend on the partition.

\textbf{Proposition 1 (finite classical compatibility).}
\(S_vS_u=S_{u+v}\) for every positive \(u,v\).

\emph{Proof.} The composed position increment is
\(up/M+Fu^2/(2M)+v(p+Fu)/M+Fv^2/(2M) =(u+v)p/M+F(u+v)^2/(2M)\); the
impulses sum to \(F(u+v)\). \(\square\)

The insertion changes the old endpoint half-kicks as well. Relative to
the old two-kick cell, the impulse changes at times \(0,u,h\) are

\[-\frac{Fv}{2},\qquad \frac{Fh}{2},\qquad -\frac{Fu}{2}.\]

Their total and their first time moment vanish. Consequently the old
endpoint position and momentum, with the convention just specified, are
preserved. Adding only the middle kick would describe a different force
history. The new chord slopes change, while retained positions and the
states between half-kicks agree.

The same result can be read directly from a discrete action. Put

\[L_h(x,y)=\frac{M(y-x)^2}{2h}+\frac{Fh}{2}(x+y).\]

\textbf{Proposition 2 (eliminating the new point).} For fixed endpoints
\(x,y\), the unique minimizer of \(L_u(x,z)+L_v(z,y)\) is

\[z_* =\frac{vx+uy}{h}-\frac{Fuv}{2M},\]

and

\[\min_z\{L_u(x,z)+L_v(z,y)\}
=L_h(x,y)-\frac{F^2uvh}{8M}. \tag{1}\]

\emph{Proof.} With \(\bar z=(vx+uy)/h\) and \(w=z-\bar z\), completing
the square gives

\[L_u(x,z)+L_v(z,y)-L_h(x,y)
=\frac{Mh}{2uv}w^2+\frac{Fh}{2}w.\]

The minimum is at \(w=-Fuv/(2M)\) and equals the displayed constant.
\(\square\)

That constant is independent of the endpoints, so its derivatives do not
change the classical endpoint momenta. It does change the action. Define

\[\ell_h(x,y)=L_h(x,y)-\frac{F^2h^3}{24M}.\]

Since \(h^3-u^3-v^3=3uvh\), equation (1) becomes the exact law

\[\min_z\{\ell_u(x,z)+\ell_v(z,y)\}=\ell_h(x,y). \tag{2}\]

This constructs the corrected endpoint action from finite insertion
compatibility. Two different orders of inserting several points agree:
the correction removed is always \(F^2(h^3-\sum_i h_i^3)/(24M)\).
Equivalently, the two-step insertion correction satisfies the
associativity identity \(d(u,v)+d(u+v,w)=d(v,w)+d(u,v+w)\).

Among continuous corrections depending only on the cell duration, the
cubic correction is determined up to \(E_{\rm ref}h\): subtracting two
solutions of the insertion equation gives the additive Cauchy equation.
This remaining freedom changes the zero of energy, leaving spectral
differences unchanged. Thus insertion consistency fixes a nontrivial
part of the action while retaining the expected energy-normalization
freedom.

For the user's example, if \(t_{2.6}\) lies at \(3/5\) of the old cell,
then \(u=3h/5\), \(v=2h/5\). The insertion removes \(3F^2h^3/(100M)\),
or \(18/25\) of that cell's original cubic correction; \(7/25\) remains.
These are exact fractions.

The geometric interpretation stays with Galileo's inertial line and
parabola. At horizontal speed \(V\), the triangle between the old chord
and the two new chords has area \(VFuvh/(4M)\) for \(F,V>0\);
multiplying by \(F/(2V)\) gives (1)'s action correction. It is the
chord-segment diagnostic of that comparison, not a Kepler swept sector.

\hypertarget{quantum-composition-and-the-arbitrary-partition-limit}{%
\subsection{3. Quantum composition, and the arbitrary-partition
limit}\label{quantum-composition-and-the-arbitrary-partition-limit}}

Supply canonical quantum kinematics \([q,p]=i\hbar\), \(\hbar>0\), and
use the unitary cell operator

\[Q_h=e^{iFhq/(2\hbar)}e^{-ihp^2/(2M\hbar)}e^{iFhq/(2\hbar)}.\]

Its coordinate kernel has phase \(L_h/\hbar\) and free normalization
\((M/(2\pi i\hbar h))^{1/2}\), with the usual positive-time branch.
Completing the same square as in Proposition 2 in the oscillatory
Gaussian integral gives

\[Q_vQ_u=e^{-iF^2uvh/(8M\hbar)}Q_h.\]

The Gaussian normalization composes too: the intermediate integral
contributes \((2\pi i\hbar uv/(Mh))^{1/2}\). Hence

\[U_h=e^{-iF^2h^3/(24M\hbar)}Q_h,
\qquad U_vU_u=U_{u+v}. \tag{3}\]

This is the constant-force propagator, with generator \(p^2/(2M)-Fq\).
The scalar correction is measurable as a relative phase when histories
are coherently compared. It cancels from the body-only channel
\(\rho\mapsto Q_h\rho Q_h^\dagger\).

For \(F\ne0\) this mechanical Hamiltonian on the line has spectrum
\(\mathbb R\): unitary translations of \(q\) shift it by arbitrary real
constants. The present limit is a unitary evolution over a finite
experiment. Its action-cost question concerns records of that
experiment; the field-theory vacuum spectral gap is a further kind of
statement.

The following finite-to-limit statement also covers situations with an
error instead of an exact scalar correction. It is an elementary
contractive form of a noncommutative sewing argument
(\href{https://arxiv.org/abs/0706.0202}{Feyel--de La
Pradelle--Mokobodzki 2007}, abstract; the stronger telescoping
hypothesis used here is stated and proved below).

\textbf{Theorem 3 (summable insertion defects).} Let \(Q(t,s)\) be
bounded operators on a fixed Banach space, \(s<t\) in a finite interval,
with \(\|Q(t,s)\|\le1\). Suppose for every \(s<r<t\),

\[\|Q(t,s)-Q(t,r)Q(r,s)\|
\le C\bigl((t-s)^p-(r-s)^p-(t-r)^p\bigr),\qquad p>1. \tag{4}\]

For a finite partition \(\pi\) of \([s,t]\), let \(Q_\pi\) be the
chronological product. There exists a unique limit \(U(t,s)\) as the
mesh tends to zero, independent of nesting, unequal cell lengths or
insertion order, and

\[\|Q_\pi-U(t,s)\|\le C\sum_{I\in\pi}|I|^p
\le C(t-s)|\pi|^{p-1}. \tag{5}\]

The limit is contractive and satisfies \(U(t,r)U(r,s)=U(t,s)\). If
\(Q(t,s)\) tends strongly to the identity as \(t\downarrow s\), so does
\(U(t,s)\).

\emph{Proof.} One insertion changes the full product by at most (4),
since all operators before and after that cell are contractions.
Successive insertions telescope the potential \(\sum_I|I|^p\), so for a
refinement \(\pi'\),

\[\|Q_\pi-Q_{\pi'}\|
\le C\left(\sum_{I\in\pi}|I|^p-
\sum_{J\in\pi'}|J|^p\right).\]

Two arbitrary partitions have a common refinement, their union. Their
products differ by at most \(C(t-s)\) times the sum of their meshes to
the power \(p-1\). Completeness gives the limit; refining a fixed
\(\pi\) proves (5). Joining partitions on \([s,r]\) and \([r,t]\) proves
composition. The one-cell bound \(\|U(t,s)-Q(t,s)\|\le C(t-s)^p\) proves
the final assertion. \(\square\)

For constant force, \(|e^{ix}-1|\le|x|\) gives (4) with \(p=3\) and
\(C=F^2/(24M\hbar)\). Thus, for duration \(T\),

\[\|Q_\pi-U_T\|
\le\frac{F^2T}{24M\hbar}|\pi|^2. \tag{6}\]

This is a constructed arbitrary-partition limit. The analogous estimate
for a general nonlinear force needs a suitable domain and stability
norm; unbounded commutators need not satisfy (4) in operator norm. The
constant-force proof does not silently supply that extension.

Neither (2) nor (3) determines a universal action unit. The classical
composition is exact, and the quantum composition works at every
supplied \(\hbar>0\). For the Newton necessity goal, the new question is
which independently justified physical record or optical law selects the
quantum composition structure and a common positive normalization.
Refinement of unobserved intermediate variables already works at fixed
\(\hbar\); an action floor concerns the physical records, not a failure
of (6).

\hypertarget{b.-harmonic-polygons-a-finite-refinement-bound-from-a-different-norm}{%
\subsection{3b. Harmonic polygons: a finite refinement bound from a
different
norm}\label{b.-harmonic-polygons-a-finite-refinement-bound-from-a-different-norm}}

The sister repository's
\href{https://github.com/arivero/newtonlean/blob/381d5e9/research/PROP_I_REALIZATION.md}{Proposition
I realization graph} and
\href{https://github.com/arivero/newtonlean/blob/381d5e9/NewtonLimitDynamics/Polygon/HarmonicStability.lean}{HarmonicStability}
(local full-read at commit \texttt{381d5e9}) separate finite stability
from convergence of force-generated polygons. The equal-cell quadratic
invariant is proved there; a norm inequality in the represented-fraction
arithmetic is a pending formalization step. The following written
argument uses a weighted sum of absolute values instead. It has not been
formalized in that repository and imports no trajectory or ODE existence
theorem.

Use its \textbf{end-kick} convention, distinct from §2's half-kicks. In
one coordinate, acceleration is \(a(x)=-wx\), with \(w\ge0\), and a cell
of duration \(h\) first drifts and then applies the force at the new
position: \[A_h\binom{x}{v}
=\begin{pmatrix}1&h\\-wh&1-wh^2\end{pmatrix}\binom{x}{v}.
\tag{6a}\] For several coordinates apply the same map to each coordinate
pair. Choose \(\rho>0\) with \(w\le\rho^2\) and define
\(\|z\|_\rho=\rho\sum_i|x_i|+\sum_i|v_i|\). Let
\(A_\pi=A_{h_n}\cdots A_{h_1}\) for a finite positive partition of
\([0,\tau]\), and assume \(\rho|\pi|\le1\).

\textbf{Proposition 3b (finite harmonic refinement control).} Put
\(B_\tau=e^{\rho\tau}\). Every such partition satisfies
\(\|A_\pi\|_\rho\le B_\tau\). If \(\pi'\) refines \(\pi\), then
\[\|(A_{\pi'}-A_\pi)z\|_\rho
\le wB_\tau\left(\sum_{I\in\pi}|I|^2
                    -\sum_{J\in\pi'}|J|^2\right)\|z\|_\rho.
\tag{6b}\] For any two admissible partitions of the same interval,
\[\|(A_\pi-A_\sigma)z\|_\rho
\le w\tau B_\tau(|\pi|+|\sigma|)\|z\|_\rho.
\tag{6c}\] If \(\rho\tau<1\), each occurrence of \(B_\tau\) in these
bounds may instead be replaced by the rational expression
\((1-\rho\tau)^{-1}\). Thus on that window the proof uses only finite
arithmetic, order and absolute-value inequalities when the data are
rational.

\emph{Proof.} In scaled coordinates \((\rho x,v)\) the absolute column
sums of \(A_h\) are \(1+wh/\rho\) and \(\rho h+|1-wh^2|\). Since
\(wh^2\le1\), both are at most \(1+\rho h\). The triangle inequality
therefore proves \(\|A_h\|_\rho\le1+\rho h\). Multiplication gives
\(\|A_\pi\|_\rho\le\prod_j(1+\rho h_j)\le e^{\rho\tau}\). For
\(\rho\tau<1\), expand the finite product: its degree-\(r\) elementary
symmetric coefficient is at most \((\rho\tau)^r\). The finite geometric
sum is at most \((1-\rho\tau)^{-1}\). This alternative bound applies to
every subproduct too.

For a split \(h+k\), direct multiplication gives \[\begin{aligned}
\Delta x&=-whk(x+hv),\\
\Delta v&=whk\,v+w^2hk^2(x+hv),
\qquad\Delta=(A_kA_h-A_{h+k})z.
\end{aligned}\tag{6d}\] The absolute column sums of the scaled
difference matrix are \[whk(1+wk/\rho),\qquad whk(1+\rho h+whk).\] They
are at most \(2whk\): use \(w\le\rho^2\) and \(\rho h+\rho k\le1\),
which also gives \(\rho h+whk\le\rho h+\rho^2hk\le1\). Thus
\(\|A_kA_h-A_{h+k}\|_\rho\le2whk\) \(=w[(h+k)^2-h^2-k^2]\). Multiplying
by all the earlier and later cell maps costs at most \(B_\tau\); their
combined duration is at most \(\tau\). Successive insertions telescope
the sum of squared cell lengths, proving (6b). For arbitrary
\(\pi,\sigma\) use their finite union as a common refinement and
\(\sum h_j^2\le\tau|\pi|\), proving (6c). \(\square\)

The estimate also controls the actual polygon between vertices. At a
time \(t\) inside a cell, let \(u\) be its elapsed duration and let
\(\pi_t\) consist of the completed cells followed by \(u\) (omit a zero
cell). The comparison state \(z_\pi^\dagger(t)=A_{\pi_t}z_0\) has
exactly the polygon position; its velocity includes a partial kick that
the actual polygon has not yet received. For the actual state
\(z_\pi(t)\), take the outgoing velocity at a completed vertex. Then
\[\|z_\pi^\dagger(t)-z_\pi(t)\|_\rho
=wu\sum_i|x_{\pi,i}(t)|
\le\frac{w}{\rho}|\pi|B_\tau\|z_0\|_\rho.
\tag{6e}\] Apply (6c) to the two clipped partitions of \([0,t]\), then
(6e) twice: \[\sup_{0\le t\le\tau}\|z_\pi(t)-z_\sigma(t)\|_\rho
\le wB_\tau(\tau+\rho^{-1})(|\pi|+|\sigma|)\|z_0\|_\rho.
\tag{6f}\] This is a finite uniform Cauchy estimate, including unequal
schedules and retained velocity. Completing the state space would
construct a limit; identifying its force law and its historical
geometric enclosure remain additional steps. No curve was supplied to
prove (6b)--(6f). The norm change supplies a concrete finite input to
the graph's P3 obligation, rather than deducing P3 from an area
invariant.

There is also a finite input to force identification. Define the
trapezoidal position sum, with no integral presumed,
\[T_\pi(t)=\sum_{j\text{ complete}}\frac{h_j}{2}(x_{j-1}+x_j)
             +\frac u2(x_{\rm last}+x_\pi(t)).\] The same end-kick
construction gives the exact identity \[v_\pi^\dagger(t)-v_0+wT_\pi(t)
=-\frac w2\left[\sum_{j\text{ complete}}h_j^2v_{j-1}
                       +u^2v_{\rm last}\right].\tag{6g}\] Indeed
\(x_j-x_{j-1}=h_jv_{j-1}\) and the kick is \(-wh_jx_j\); subtracting a
trapezoid from \(h_jx_j\) leaves \(h_j^2v_{j-1}/2\). The partial kick
supplies the identical calculation with \(u\). Each departure velocity
has absolute sum at most \(B_\tau\|z_0\|_\rho\), and
\(\sum h_j^2+u^2\le t|\pi|\). For the actual velocity add the missing
partial kick \(wu x_\pi(t)\) from (6e), obtaining
\[\|v_\pi(t)-v_0+wT_\pi(t)\|_1
\le wB_\tau(t/2+\rho^{-1})|\pi|\|z_0\|_\rho.
\tag{6h}\] The position balance
\(x_\pi(t)-x_0=\sum_{j\text{ complete}}h_jv_{j-1}+uv_{\rm last}\) is
exact. At \(t=0\) all sums are empty; at a completed vertex \(u=0\), so
the identities retain the declared outgoing velocity. Thus the candidate
curve family is paired with a vanishing finite force-balance residual,
rather than only a stability estimate. Passing these sums to a
continuous force law and proving the relevant geometric enclosure remain
separate obligations.

\textbf{Astra's graph check (2026-09-29).} The printed-stage edges from
Lemma III Cor. 4 to Proposition I record a textual citation. They do not
supply convergence of the generated polygon family: the lemma's
given-curve approximation and the graph's P3 realization have different
inputs. The graph's attached
\href{https://github.com/arivero/newtonlean/blob/381d5e9/NewtonLimitDynamics/Polygon/Enclosure.lean}{sector\_ratio\_reconstruction}
(local full-read) explicitly assumes lower and upper limits and a
geometric enclosure. Having that formal reference does not discharge
those hypotheses. This check concerns graph metadata and code; the
printed passages retain their existing repository attribution. Astra
checked (6b)--(6h), including the single \(B_\tau\) factor and the
partial-cell sign in (6g). The finite estimate supplies P3's Cauchy
input and (6h) supplies force-consistency data; neither silently closes
the graph's limiting geometric or force-identification nodes. This is an
internal mathematical check, separate from Lean compilation.

The user's alternative (2026-09-29) is to obtain surviving quantities by
repeated block renormalization, without presupposing a limiting physical
curve. In this fixed harmonic model the Cauchy bound forces a limit in
the usual real-state completion. It does not establish that the full
body--apparatus state has the same completion when refinement adds bath
and memory variables. For that construction, specify the retained
observables and their blocking maps; §7's consistency estimate is a
distinct route to their limiting laws. A finite area identity alone
establishes neither kind of limit.

Here \(w\) has units of inverse time squared and \(\rho\) of inverse
time; the norm has velocity units, and every coefficient multiplying it
in (6b)--(6f) is dimensionless. At \(w=0\) the insertion defect and all
partition discrepancies vanish exactly. For fixed \(\tau,w,\rho\) the
bounds tend to zero with both meshes, independently of their nesting.
Mass enters only through \(w\) if a spring force is written \(-mw x\).
No action normalization appears. To use this mechanical bound for a
record, its discrepancy must still vanish after whitening by the
complete conditional record covariance, as tested in
\href{sed-closure-under-recording.md\#63-ideas-supplied-by-the-current-lean-formalization}{the
recording note, §6.3}.

\hypertarget{c.-all-refinement-sequences-and-two-microscopic-kick-orders}{%
\subsection{3c. All refinement sequences and two microscopic kick
orders}\label{c.-all-refinement-sequences-and-two-microscopic-kick-orders}}

The user's proposed criterion is independence from \textbf{every
admitted scaling sequence}, including unequal and nonnested partitions.
The local dynamics and physical initial data must be specified
throughout. One repeated dyadic sequence alone would not establish that
criterion. Here the finite estimate also permits a change of kick order
in every cell, and constructs the harmonic curve and its force law.

Retain \(w\ge0\), \(\rho>0\), \(w\le\rho^2\) and the norm of §3b. Let
\(D_h(x,v)=(x+hv,v)\) and \(K_h(x,v)=(x,v-whx)\). The end-kick and
start-kick cell maps are respectively
\[A_h=K_hD_h=\begin{pmatrix}1&h\\-wh&1-wh^2\end{pmatrix},
\qquad C_h=D_hK_h=\begin{pmatrix}1-wh^2&h\\-wh&1\end{pmatrix}.
\tag{6i}\] For each cell choose either map, independently of the choices
in other cells. Write \(T_\pi\) for their chronological product. Require
\(\rho|\pi|\le1\), which every vanishing-mesh sequence eventually
satisfies. Initial position and velocity \(z_0\) and \(w\) stay fixed.

\textbf{Proposition 3c (harmonic scheme universality).} With
\(B_\tau=e^{\rho\tau}\), \[\begin{aligned}
\|T_\pi\|_\rho&\le B_\tau,\\
\|(T_\pi-A_\pi)z_0\|_\rho
 &\le wB_\tau\sum_i h_i^2\|z_0\|_\rho,\\
\|(T_\pi-T_\sigma)z_0\|_\rho
 &\le2w\tau B_\tau(|\pi|+|\sigma|)\|z_0\|_\rho.
\end{aligned}\tag{6j}\] No nesting or common kick choices are required.
Linearly interpolate both vertex state components to obtain a continuous
displayed path \(\overline z_\pi\). Then
\[\sup_{t\le\tau}\|\overline z_\pi(t)-\overline z_\sigma(t)\|_\rho
\le2wB_\tau(\tau+\rho^{-1})(|\pi|+|\sigma|)\|z_0\|_\rho.
\tag{6k}\]

\emph{Proof.} In scaled coordinates the absolute column sums of \(C_h\)
are \(|1-wh^2|+wh/\rho\) and \(1+\rho h\). Under the stated mesh bound
they are at most \(1+\rho h\), as are those of \(A_h\). Every mixed
product consequently has norm at most \(B_\tau\). Moreover
\[A_h-C_h=wh^2\operatorname{diag}(1,-1),\qquad
\|A_h-C_h\|_\rho=wh^2.\] Replace one factor at a time. Its prefix and
suffix have disjoint durations whose sum is at most \(\tau\), so their
combined amplification is at most one \(B_\tau\). This proves the second
line of (6j). Combine it for both partitions with (6c) to prove the
third.

Interpolation between corresponding vertices preserves the
same-partition mixed-versus-end bound. The displayed velocity of an
end-kick polygon differs from its actual velocity by at most the cell
kick, hence by \((w/\rho)B_\tau|\pi|\|z_0\|_\rho\). Use these two
interpolation corrections and (6f), then the two mixed-versus-end
comparisons, to obtain (6k). \(\square\)

The vertex state is the output of the preceding cell and the input to
the next. A next-cell start kick occurs immediately after that state. At
an end-kick/start-kick interface there are two impulses; the vertex
state is between them. Each adjacent velocity differs from the displayed
vertex velocity by at most one cell-kick bound. At \(t=0\), the
displayed initial velocity is \(v_0\); a first start kick changes the
outgoing velocity by a quantity vanishing with the mesh. On either kind
of cell the actual position equals its displayed linear interpolation,
and the actual velocity differs by at most
\((w/\rho)B_\tau|\pi|\|z_0\|_\rho\).

Completeness of real continuous path space constructs a limit from (6k),
initially on any one vanishing-mesh sequence. The same bound then gives
that identical limit for every admitted partition sequence and every
sequence of kick choices. Actual velocity representatives converge
uniformly too. No limiting curve was used in these estimates.

To identify its force law, first choose an all-end-kick sequence. Its
finite position balance and (6h) pass through uniform convergence to
\[x(t)=x_0+\int_0^t v(s)\,ds,\qquad
v(t)=v_0-w\int_0^t x(s)\,ds.\tag{6l}\] The trapezoidal sum in (6h) is
exactly the integral of the finite polygon position. The position
balance is its finite drift integral. Uniform convergence passes these
ordinary integrals to the limit. Thus \(x'=v\) and \(v'=-wx\). Every
mixed construction has this same limit by (6k). This proves trajectory
existence for the specified harmonic model; the printed Newton enclosure
and the sister Lean formalization of P3 remain separate. At \(w=0\),
every discrepancy vanishes exactly. The estimates have the units and
norm of §3b.

The RG comparison concerns universality of retained descriptions. The
finite block map is the full product of its daughter matrices;
regrouping that product associates exactly. Replacing it by one bare
cell map has a defect, which the estimates control. Thus the curve is
obtained from the whole compatible refinement family. Its existence does
not require exact equality of the two finite conventions.

The source connection is explicit but has additional hypotheses. In
\href{https://arxiv.org/pdf/1405.1548v3}{'t Hooft's automaton book,
§§5.4 and 20.8} {[}@tHooft2015CAI{]} (passages read), he postulates an
ontological basis and invokes RG to connect microscopic automaton
Hamiltonians to large-scale field theory. Our inference is the
universality question: which structures survive a change of finite
description? His deterministic ontology is an additional premise. His
own
\href{https://webspace.science.uu.nl/~hooft101/gthpub.html}{publication
list} (author metadata read) dates the renormalization papers to
1971--73 and deterministic-QM papers to 1988--90; an immediate causal
transition is not established by those dates.

\textbf{Mathematical review (GPT-6 Astra, 2026-09-30): ACCEPT.} The
single \(B_\tau\) amplification, interpolation and endpoint convention,
completion and force-law passage were checked. The 't Hooft source
attribution is a separate passage check by Sol. The result concerns
fixed harmonic dynamics, not arbitrary microscopic laws or an action
floor.

\hypertarget{d.-scale-invariance-cancelled-masses-and-the-action-that-remains}{%
\subsection{3d. Scale invariance, cancelled masses and the action that
remains}\label{d.-scale-invariance-cancelled-masses-and-the-action-that-remains}}

The user's scale comparison adds an exact obstruction to identifying an
action unit from curve compatibility alone. An overall action
normalization can disappear from every classical trajectory equation and
remain in canonical momentum, energy and recorded phase.

\textbf{Proposition 3d (classical normalization under blocking).} For a
discrete action \(S_\pi\) and any dimensionless constant \(c_*>0\),
stationary configurations of \(c_*S_\pi\) are unchanged, and
\[\inf_z c_*S_\pi(x,z,y)=c_*\inf_zS_\pi(x,z,y).\tag{6m}\] The minimizing
configurations, when attained, are the same. This follows directly by
multiplying the first variation and the order inequalities by \(c_*\).
Consequently any classical insertion or blocking law built from these
configurations commutes with this rescaling. Refinement independence of
their curve fixes no overall action normalization.

For a free body, \(S=m\int|\dot q|^2dt/2\), \(m>0\); fixed initial
position and \textbf{velocity} give a trajectory independent of \(m\).
Momentum \(p=m\dot q\) and action retain it. For uniform acceleration,
put \(F=ma\) in §2's corrected cell action: \[\ell_h(x,y)=m\left[
\frac{(y-x)^2}{2h}+\frac{ah}2(x+y)-\frac{a^2h^3}{24}\right].
\tag{6n}\] The curve and point elimination are independent of the
overall \(m\); the cubic action correction is \(ma^2h^3/24\). Mass
cancellation in the motion therefore leaves the phase-space action
question intact.

For a test body in a prescribed external gravitational potential, assume
equality of its passive gravitational and inertial masses. On a
collision-free interval,
\[S=m\int\left[\frac{|\dot q|^2}2+\frac{GM}{|q|}\right]dt,
\qquad \ddot q=-GM\frac q{|q|^3}.\tag{6o}\] The test mass cancels; the
prescribed source's \(GM\) remains. For two moving gravitating bodies
the relative equation contains \(G(M+m)\), so the test-body
qualification matters. Cancellation from these equations is not a
classification as an RG-irrelevant coupling.

Under \(q_\lambda(t)=\lambda q(t/\theta)\) and a correspondingly
rescaled interval, free action scales by \(\lambda^2/\theta\) at fixed
\(m\). Gravity with fixed \(GM\) requires \(\theta=\lambda^{3/2}\) and
its action scales by \(\lambda^{1/2}\). Harmonic covariance sends \(w\)
to \(w/\theta^2\). These parameter transformations are different from
taking the partition mesh to zero at fixed physical \(w\), \(GM\) and
initial data.

In four-dimensional pure Yang--Mills, use the connection convention
\(F=dA+A\wedge A\). Give all four coordinates length units, with
\(x_4=ct_E\), and write
\[\frac{S_E}{\hbar}=\frac1{4g^2}\int F^a_{\mu\nu}F^a_{\mu\nu}
\,d^4x.\] The classical dilation
\(A_\lambda(x)=\lambda^{-1}A(x/\lambda)\) gives
\(F_\lambda(x)=\lambda^{-2}F(x/\lambda)\); its fourth-power factor
cancels the volume factor. The coupling is dimensionless, and its
overall \(g^{-2}\) factor also cancels from the classical vacuum
equations. Quantum running makes this normalization relevant to
fluctuations. The gauge comparison already supplies \(\hbar>0\); it does
not derive that assumption from Newton's trajectory.

With inverse-length renormalization scale \(\mu\), the standard
weak-coupling beta function for pure \(SU(3)\) is
\[\mu\frac{dg}{d\mu}=-b_0g^3+O(g^5),\qquad
b_0=\frac{11}{16\pi^2}.\] The normalization is checked in
\href{https://arxiv.org/pdf/1006.4518v3}{Lüscher, equation (2.31)}
{[}@Luscher2010WilsonFlow{]} (passage read, \(N=3,N_f=0\); §3.3 also
describes quantum breaking of classical scale invariance). In the
\textbf{one-loop truncated flow}, dimensional transmutation gives
\[\Lambda_1=\mu e^{-1/(2b_0g^2(\mu))},\qquad
g^{-2}(a_n)=g^{-2}(a_0)+2b_0\log(a_0/a_n),\quad\mu=1/a.
\tag{6p}\] Differentiating \(g^{-2}\) gives \(2b_0\) with respect to
\(\log\mu\), so \(\Lambda_1\) is invariant and consecutive logarithmic
increments telescope. Equation (6p) holds along every scaling sequence
in the positive-coupling domain \(a\Lambda_1<1\); every sequence
\(a_n\to0\) eventually enters it. Monotonicity and fixed scaling ratios
are unnecessary within that domain. Higher-loop running and scheme
normalization are separate from this explicit truncated formula.

\textbf{The overall factor in a retained weight.} For a dimensionless
fixed reference measure \(\nu\) and a finite positive integral, define
Euclidean or Gibbs elimination by \[\mathcal B_\hbar(S)(x)
 =-\hbar\log\int e^{-S(x,z)/\hbar}\,d\nu(z),\qquad \hbar>0.\] Then the
user's denominator observation remains exact under blocking:
\[e^{ic_*S/\hbar}=e^{iS/(\hbar/c_*)},\qquad
\mathcal B_\hbar(c_*S)=c_*\mathcal B_{\hbar/c_*}(S).
\tag{6q}\] The first identity concerns real-time phase; the second
concerns a positive Euclidean weight. Its proof is substitution into the
logarithm. Fixed normalizations independent of \(x\) change only an
additive action constant. Compatible product measures make successive
eliminations associate by Fubini. A reference measure that also changes
with \(c_*\) must be compared separately.

Once the blocking maps and measures are specified, marginalization
therefore determines the retained weight at each resolution.
Reexpressing these actions at a common cutoff defines an RG
transformation. It can generate interactions absent from the starting
action, remain at a fixed point, or change only by normalization.
Neither a differentiable flow nor closure in finitely many couplings
follows from the identity. The higher-dimensional gauge interactions in
§6 give its concrete use. Interpreting a positive weight as quantum
theory also requires the supplied reconstruction hypotheses; the
identity alone is not a derivation of quantum mechanics.

At fixed \(\hbar\), multiplying \(S\) changes fluctuation weights
exactly as replacing its action denominator by \(\hbar/c_*\). Classical
stationarity cannot distinguish those weights. In the pure-gauge
connection convention, put \(I[A]=\frac14\int F^2d^4x\) and
\(S_E=A_*I[A]\), with \(A_*\) in action units. Then \(g^2=\hbar/A_*\).
Changing \(A_*\) to \(c_*A_*\) at the same reference \(\mu\) gives,
within the one-loop truncated flow, \[g_{c_*}^2=\frac{g^2}{c_*},\qquad
\Lambda_{1,c_*}=\mu e^{-c_*/(2b_0g^2)}
              =\mu(\Lambda_1/\mu)^{c_*}.\tag{6r}\] Thus an overall
action factor can change a generated inverse length exponentially, even
though it cancels from the classical vacuum equations. It is not a
linear change of units for \(\Lambda_1\).

Keeping constants explicit gives a parallel mechanical calculation. For
\(L,T,m,GM>0\), write \(q=Ly\), \(t=Ts\) in (6o). Then
\[\frac{S}{\hbar}=\frac{mL^2}{\hbar T}
 \int\left[\frac{|y'|^2}2+
                  \frac{GMT^2/L^3}{|y|}\right]ds.\] The dimensionless
motion parameter \(GMT^2/L^3\) and action prefactor \(mL^2/(\hbar T)\)
have different roles. Choosing \(T=\sqrt{L^3/(GM)}\) sets the first to
one and leaves \(m\sqrt{GML}/\hbar\) in the action. At the conversion
length \(r_g=GM/c^2\), that prefactor is \(\alpha_G=GMm/(\hbar c)\).
This choice introduces no relativistic dynamics into the Newtonian
model. For a free body only the kinetic term and its action prefactor
remain. Use \(h_P=2\pi\hbar\) for Planck's constant; the time-cell
duration \(h\) in (6n) is a different quantity.

This shows consistency of running along arbitrary admitted scale
sequences. It supplies no bound on the full blocked action, continuum
axioms or physical spectrum. The remaining spectral requirement is
\[0<\frac{E_{\rm gap}}{\hbar c\Lambda}<\infty,\] with surviving
finite-energy spectral weight, as in
\href{three-continuum-limits.md}{the continuum comparison}. Thus the
project tracks three distinct survivors: a motion or observable law
independent of refinement, a physical dimensional scale, and action
normalization. The harmonic curve construction proves the first in its
model. Neither it nor mass cancellation selects a positive action unit,
and dimensional transmutation alone proves no mass gap.

\textbf{Scaling check (GPT-6 Astra, 2026-09-30): ACCEPT at the stated
classical and one-loop levels.} The passive-mass scope, dilation
exponents, running sign, positive-coupling domain, retained-weight
identity and explicit \(c,\hbar,G\) dimensions were checked.

\hypertarget{gauge-refinement-geometry-marginalization-and-reflection}{%
\subsection{4. Gauge refinement: geometry, marginalization and
reflection}\label{gauge-refinement-geometry-marginalization-and-reflection}}

Take a finite oriented lattice \(\Gamma\) and a refinement \(\Gamma'\).
For an old edge represented by a path of fine edges, define

\[p_{\Gamma'\Gamma}(U)_e=U_{e_1}\cdots U_{e_k},\]

using inverses when an orientation reverses. These maps compose exactly
under further refinements and are gauge covariant: transformations at
the new interior vertices cancel in the product. Holonomies of old loops
are therefore represented exactly on the fine lattice.

For a probability measure the consistency equation is stronger:

\[p_{\Gamma'\Gamma*}\mu_{\Gamma'}=\mu_\Gamma. \tag{7}\]

At finite volume the left side can always be formed. Its effective
Boltzmann weight is the integral of the fine weight over the fibres of
\(p\), with the Haar disintegration. This defines an effective action up
to a normalization constant. Equation (7) with a proposed coarse action
is a theorem to prove, not a consequence of having multiplied the links
correctly.

Uniformly halving a \(D\)-cube produces \(2^D\) subcubes: four, eight or
sixteen for \(D=2,3,4\). Local subdivisions and anisotropic spacetime
refinements are possible as well, provided neighbouring cells and the
projection maps remain compatible. The uniform grid is a convenient
sequence through a larger refinement problem.

\textbf{Time refinement alone versus a field continuum.} At fixed
spatial spacing \(a_s\), a Hamiltonian lattice theory has the exact
temporal law \(T_{a_s}(t)=e^{-tH_{a_s}/\hbar}\) and
\(T_{a_s}(u+v)=T_{a_s}(v)T_{a_s}(u)\). Taking the time step to zero
leaves the spatial cutoff in place. Refining space adds degrees of
freedom and changes the Hilbert space, so Theorem 3 on one fixed space
cannot be applied without comparison maps. An isotropic Euclidean
refinement removes both cutoffs together; an anisotropic route must also
match electric and magnetic normalizations to recover the same physical
speed \(c\). This is the additional task in the four-cube comparison
beyond inserting a time in mechanics.

One essential property can be carried exactly. Suppose a fine measure is
reflection positive, \(p\) commutes with time reflection \(\Theta\), and
every coarse positive-time observable pulls back to a fine positive-time
observable. Then its pushforward is reflection positive:

\[\int\overline{\Theta f}\,f\,d(p_*\mu)
=\int\overline{\Theta(f\circ p)}(f\circ p)\,d\mu\ge0. \tag{8}\]

The support condition on \(p\) matters; a block straddling the
reflection plane cannot be assumed to have it. Equation (8) supplies
positivity for an exact blocking map with those properties. Subsequent
truncation of the effective action must preserve it separately.

\hypertarget{the-exactly-soluble-subdivision-law-in-11-dimensions}{%
\subsection{\texorpdfstring{5. The exactly soluble subdivision law in
\(1+1\)
dimensions}{5. The exactly soluble subdivision law in 1+1 dimensions}}\label{the-exactly-soluble-subdivision-law-in-11-dimensions}}

Let \(G=SU(3)\) with normalized Haar measure. Choose the group metric so
that \(-\Delta_G\chi_R=C_2(R)\chi_R\) and \(C_2(\mathbf3)=4/3\). Let
\(\lambda_2>0\) have units of inverse area and define the face kernel

\[k_A(U)=\sum_R d_R\chi_R(U)e^{-\lambda_2 A C_2(R)/2}.\]

It is the heat kernel at group heat-time \(\lambda_2A\). Character
orthogonality gives the exact convolution identity

\[\int_G k_{A_1}(XU^{-1})k_{A_2}(UY)\,dU
=k_{A_1+A_2}(XY). \tag{9}\]

Indeed, integrating two characters contributes
\(\delta_{RS}\chi_R(XY)/d_R\), so their heat exponents add and one
factor \(d_R\) remains. When an interior edge splits a face, (9)
integrates out precisely that new edge. Splitting an edge alone uses
invariance of Haar measure under multiplication. These elementary moves
yield (7) for the heat-kernel surface measure with density
\(Z^{-1}\prod_f k_{A_f}(U_{\partial f})\).

This is the established lattice-to-continuum construction of
two-dimensional Yang--Mills. Lévy proves subdivision consistency and
constructs continuum random holonomy, including the continuity needed to
extend beyond a nested set of graph edges
(\href{https://arxiv.org/pdf/math/0101239}{Lévy, \emph{Yang--Mills
Measure on Compact Surfaces}, Theorem 1.6.1 and §§2.5--2.10}, passage).

Here the counterpart of Newton's unequal-time insertion is division into
any two positive areas, \(A=A_1+A_2\). The law already knows how to
remove the new variable. It also keeps its coupling parameter:
composition works for every \(\lambda_2>0\).

\textbf{The spectrum must still be identified.} On a spatial circle of
circumference \(L\), a Euclidean duration \(t\) has area \(Lct\). The
physical Hilbert space is \(L^2(G)^G\), the class functions, and the
character expansion gives

\[E_R=\frac{\hbar c\lambda_2L}{2}C_2(R),\qquad
E_1-E_0=\frac23\hbar c\lambda_2L\quad\text{for }SU(3). \tag{10}\]

These are global electric-flux energies. Pure \(1+1\)-dimensional
Yang--Mills has no propagating transverse gluon, and (10) grows with
\(L\). Thus the solved refinement limit supplies neither a finite
infinite-line glueball mass nor the four-dimensional spectrum. Its value
for this programme is the exact composition mechanism and a spectral
problem whose operator and boundary conditions are explicit.

There is a dimensional mass unit \(\hbar\sqrt{\lambda_2}/c\) and a
string tension unit \(\hbar c\lambda_2\). The absence of local particle
excitations has a dynamical cause. This corrects the older dimensional
claim in \href{low-dimensional-mass-gap.md}{G07} that no mass unit can
be formed in two dimensions.

\textbf{Other \(1+1\) gap theories sharpen the comparison.} The
Schwinger model, \(U(1)\) gauge theory with one massless charged Dirac
fermion, has an exactly soluble massive local boson
(\href{https://doi.org/10.1103/PhysRev.128.2425}{Schwinger 1962},
abstract). Its axial singlet symmetry is anomalous. With several
massless flavours, the spectrum instead includes a massive boson and
\(N_f-1\) massless bosons
(\href{https://arxiv.org/pdf/1508.01685}{Keegan 2015, §3}, passage).
Coleman's theorem also constrains continuous symmetry breaking under its
relativistic \(1+1\) hypotheses
(\href{https://doi.org/10.1007/BF01646487}{Coleman 1973}, metadata).
Thus lower dimension supplies exact theories with real spectral content,
but the four-dimensional pion's symmetry-breaking mechanism must be
checked afresh there. These examples serve the comparison; they do not
enlarge the main proof goals.

\hypertarget{what-changes-in-12-and-13-dimensions}{%
\subsection{\texorpdfstring{6. What changes in \(1+2\) and \(1+3\)
dimensions}{6. What changes in 1+2 and 1+3 dimensions}}\label{what-changes-in-12-and-13-dimensions}}

An interior link of a hypercubic \(D\)-dimensional lattice borders
\(2(D-1)\) plaquettes. In \(D=2\) the two-face convolution closes. In
\(D=3\) or \(4\), integrating a shared link couples four or six incident
plaquettes. The result generally depends on several boundary loops,
rather than on one coarse plaquette with one adjusted coefficient.

There is an exact local representation of this operation. Expand the
face weights in characters. For chosen incident representations, the
link integral is

\[P_{\rm inv}=\int_G\bigotimes_i R_i^{\sigma_i}(U)\,dU, \tag{11}\]

where a reversed orientation uses the dual representation. Haar
invariance and unitarity give \(P_{\rm inv}^*=P_{\rm inv}\) and
\(P_{\rm inv}^2=P_{\rm inv}\), with \(\|P_{\rm inv}\|\le1\). It is the
orthogonal projector onto invariant tensors. In two dimensions Schur
orthogonality gives (9); in higher dimensions the additional
representation labels and invariant tensors carry the boundary
interaction. Formula (11) is an exact finite integration identity, not a
physical transfer-matrix gap.

The analogy with (1) is now precise. Newton's quadratic elimination adds
an endpoint-independent scalar; after the cubic correction, the same
cell law closes. Gauge-field elimination produces an entire
boundary-dependent effective interaction. A continuum argument must
control that interaction and the observables along with the coupling.

Use a dimensionally explicit convention

\[\frac{S_E}{\hbar}=\frac1{4\lambda_D}
\int F^a_{\mu\nu}F^a_{\mu\nu}\,d^Dx,
\qquad [\lambda_D]={\rm length}^{D-4}.\]

In the classical-coupling convention of G07,
\(\lambda_D=\hbar g_{\rm cl,D}^2\). The dimensionless lattice coupling
is \(g_{\rm lat}^2=\lambda_Da^{4-D}\); in \(D=4\) it must also be run
with the cutoff.

\begin{longtable}[]{@{}
  >{\raggedright\arraybackslash}p{(\columnwidth - 4\tabcolsep) * \real{0.3333}}
  >{\raggedright\arraybackslash}p{(\columnwidth - 4\tabcolsep) * \real{0.3333}}
  >{\raggedright\arraybackslash}p{(\columnwidth - 4\tabcolsep) * \real{0.3333}}@{}}
\toprule\noalign{}
\begin{minipage}[b]{\linewidth}\raggedright
Spacetime
\end{minipage} & \begin{minipage}[b]{\linewidth}\raggedright
Refinement parameter at fixed physical coupling
\end{minipage} & \begin{minipage}[b]{\linewidth}\raggedright
Physical gap question
\end{minipage} \\
\midrule\noalign{}
\endhead
\bottomrule\noalign{}
\endlastfoot
\(1+1\) & \(\lambda_2a^2\) & Exact holonomy construction; circle flux
spectrum (10) \\
\(1+2\) & \(\lambda_3a\) & \(E_{\rm gap}=C_3\hbar c\lambda_3\), with
\(0<C_3<\infty\) to prove \\
\(1+3\) & \(g_0^2(a)\), dimensionless &
\(E_{\rm gap}=C_4\hbar c\Lambda\), with scale and positive finite
\(C_4\) to establish \\
\end{longtable}

In \(1+2\) dimensions the coupling supplies an inverse length, and the
ultraviolet problem is superrenormalizable. Dimensional analysis fixes
the unit in the table, conditional on constructing the theory; it does
not determine \(C_3\) or its sign. At finite physical \(L\), the form is
\(E(L)=\hbar c\lambda_3 f(\lambda_3L)\), with the large-\(L\) limit
still to control. The small-box constant-mode approximation is a
distinct limit from small lattice spacing at fixed box size.

\textbf{Source boundary in three dimensions.} Chandra, Chevyrev, Hairer
and Shen construct a renormalized stochastic Yang--Mills--Higgs flow on
the three-torus, locally in its auxiliary stochastic time, allowing
possible finite-time blow-up. Global survival and an invariant measure
remain open in that construction
(\href{https://arxiv.org/pdf/2201.03487}{2024 paper, introduction},
passage). It therefore does not establish a finite-volume OS quantum
theory or its physical gap. The earlier
\href{three-dimensional-gap-one-function.md}{one-function note}
overstated this input; its dimensional reduction is conditional on
existence.

The Karabali--Nair Hamiltonian approach supplies another mechanism, but
its current-variable kinetic scale must be distinguished from a proved
gap on the full physical Hilbert space
(\href{https://arxiv.org/pdf/hep-th/9602155}{Karabali--Nair 1996, eqs.
(28)--(29)}, passage). Nair's 2026 treatment presents an outline of a
proof and identifies the gap as open
(\href{https://arxiv.org/abs/2608.10133}{\emph{Towards a Proof of Mass
Gap in 3d Yang--Mills Theory}}, abstract and introductory passage). The
\(1+2\) case is a genuine target with simpler scaling, not a solved
theorem to import wholesale.

In \(1+3\) dimensions the leading coupling is marginal. Perturbatively,
one blocking step changes it by

\[g_0^{-2}(2a)=g_0^{-2}(a)-2b_0\log2+O(g_0^2(a)),
\qquad b_0=\frac{11}{16\pi^2}\quad(SU(3)).\]

The ultraviolet running sets the scale convention; the gap requires
control through the regime where that expansion no longer suffices. This
is the role of H1 in the existing
\href{mass-gap-conditional-theorem.md}{mass-gap map}, with construction
and nontriviality still explicit. The dimension table supplies a
hierarchy of constructive problems rather than an implication from one
dimension's gap to another's.

\hypertarget{what-estimate-would-construct-the-gauge-limit}{%
\subsection{7. What estimate would construct the gauge
limit?}\label{what-estimate-would-construct-the-gauge-limit}}

Here is a probability counterpart of Theorem 3. Fix a finite physical
box and nested lattices \(\Gamma_n\), \(a_n=2^{-n}a_0\), with compatible
projections \(p_{mn}\), \(m>n\). Let \(\mu_n\) be probability measures
and use total variation with convention
\(\|\mu-\nu\|_{\rm TV}=\sup_A|\mu(A)-\nu(A)|\).

\textbf{Proposition 4 (summable consistency errors).} If

\[\|p_{n+1,n*}\mu_{n+1}-\mu_n\|_{\rm TV}\le e_n,
\qquad \sum_{n=0}^\infty e_n<\infty, \tag{12}\]

then \(p_{mn*}\mu_m\) converges in total variation as \(m\to\infty\) to
a probability measure \(\nu_n\). These limits are exactly consistent,
\(p_{n+1,n*}\nu_{n+1}=\nu_n\), with

\[\|\nu_n-\mu_n\|_{\rm TV}\le\sum_{j\ge n}e_j.\]

\emph{Proof.} Pushforward contracts total variation, so the difference
of the \(m\)-th and \((m+1)\)-st projected measures is at most \(e_m\).
The series is Cauchy in the space of finite measures. Limits preserve
normalization and commute with each fixed pushforward. \(\square\)

The consistent measures define a measure on the inverse limit of these
compact graph-configuration spaces. This is a measure on generalized
holonomies for the chosen graphs. Extending it to physical continuum
fields or all relevant paths, establishing Euclidean covariance and OS
regularity, and proving nontriviality need additional estimates. A
fixed-volume bound in (12) also needs a local version uniform in the box
to support \(L\to\infty\).

For \(e_n\le C a_n^\alpha\), \(\alpha>0\), the tail is bounded by
\(Ca_n^\alpha/(1-2^{-\alpha})\). The Newton estimate uses this kind of
summability. In a fixed \(D\)-dimensional volume, a per-cell
contribution of size \(a^{D+\alpha}\) would sum to order \(a^\alpha\);
obtaining an actual bound on normalized observables or measures from
that power counting is a separate step. Marginal terms with accumulated
logarithms must first be absorbed into running couplings and
counterterms. A geometrically shrinking cell by itself supplies no
summable estimate.

Total variation in (12) is a sufficient, potentially overstrong norm. A
useful implementation can instead bound a separating class of
renormalized local observables and prove tightness in a suitable
topology. The topology, observables and constants must be named. In
higher dimensions unsmeared loop holonomies may themselves require
renormalization, so the bare graph projective limit alone is too weak a
target for the physical theory.

Finally, consistency of measures supplies no energy threshold. To retain
a positive physical gap, carry a uniform large-time bound

\[\langle O^*e^{-t(H-E_0)/\hbar}O\rangle_c
\le K_Oe^{-E_-t/\hbar},\qquad E_->0,\]

for a dense family of reconstructed local vectors, with surviving
finite-energy spectral weight. For the pion benchmark, the preserved
Ward identity and broken-symmetry residue instead require a soft
channel. The \href{three-continuum-limits.md}{spectral criteria} state
the two limiting conclusions. Short-distance consistency and
long-distance spectral control are complementary parts of one
construction.

\hypertarget{consequence-for-state-and-the-final-paper}{%
\subsection{8. Consequence for STATE and the final
paper}\label{consequence-for-state-and-the-final-paper}}

The final paper should be organized around \textbf{local insertion laws,
their composition, and the quantity carried through the limit}. The
sequence is now concrete: the Newton action elimination (1), the
arbitrary-partition construction (4)--(6), exact surface subdivision
(9), and the higher-dimensional consistency estimate (12), followed by
the physical spectral or record-cost estimate. Pion physics tests which
symmetries those estimates preserve.

The next bounded field-theory task is \textbf{one \(1+2\)-dimensional
blocking step with its boundary interaction retained}. Start with a
subdivided finite cube and compatible reflection plane; specify its
exact boundary weight via (11), then choose a renormalized local
observable norm in which the difference from a proposed coarse weight
can be bounded. The deliverable is one explicit estimate, with
dependence on \(a\), \(\lambda_3\) and boundary data, or the exact term
preventing that estimate. Its purpose is to test summability, not to fit
the answer to a one-plaquette action. The small-volume H3 valley-lifting
task remains the separate spectral mechanism to connect to this
construction.

For Newton, §§3b--3c now construct the harmonic curve and its force law,
uniformly over unequal, nonnested refinements and arbitrary mixtures of
the two force-consistent kick orders. The norm avoids the sister
formalization's pending quadratic inequality; its printed enclosure and
formal P3 remain separate. Section 3d fits the user's RG/'t Hooft and
scale-invariance threads through a normalization obstruction and an
explicit truncated running law: curve consistency, dimensional scale and
action normalization require different proofs. Exact retained weights
explain the appearance of RG as a change of resolution; useful continuum
and spectral estimates remain its test. The
\href{sed-closure-under-recording.md\#814-the-limiting-momentum-record-and-retained-controller-recoil}{physical
recording construction, §§8.11--8.14} also identifies an adaptive
body--record continuum law and smeared recoil from finite pointers, with
limiting-record closure at its supplied scale.
\href{sed-closure-under-recording.md\#815-a-nonlinear-copy-isolates-the-missing-curvature-contribution}{The
nonlinear-copy test, §8.15} shows that the same Gaussian classical
preparation fails beyond affine body copies. Its §8.16 derives the exact
quantum generator correction, which changes conditional momentum while
preserving record statistics, and excludes a positive classical
transition kernel with that generator on the full canonical phase space.
Next Newton lemma: compositional closure of nonlinear physical recording
from an independent conditioning rule; reject an imported quantum kernel
or a zero branch satisfying the same premises. The gauge frontier
remains the
\href{su2-midplane-small-field.md\#33-retained-barriers-and-nonlinear-connectors-in-a-weak-recoupling-strip}{weak-strip
marked density norm}.
\end{document}
